% photon-ring-structure.tex — Photon ring sub-structure schematic
% Inputable TikZ fragment for fig:photon-ring-structure
% Requires: tikz with calc, arrows.meta, decorations.markings (loaded by ehbook.cls)
%
% Two-panel figure.
% Left: ray trajectories for n=0 (gold), n=1 (blue), n=2 (crimson)
%   around a Schwarzschild black hole (M=1).
% Right: observer's image plane showing sub-ring annuli.
%
\begin{tikzpicture}[
    >={Stealth[length=5pt]},
    font=\footnotesize,
  ]

  % ── Panel (a): Ray trajectories ──────────────────────────────────────

  \begin{scope}[shift={(0,0)}]

    % Scale: 0.55 cm per M.  Black hole r=2M, photon sphere r=3M.
    \def\scl{0.55}

    % Black hole (filled circle, r_+ = 2M)
    \fill[black] (0,0) circle ({2*\scl});

    % Photon sphere (dashed circle, r = 3M)
    \draw[EHMarginGray, dashed, thin] (0,0) circle ({3*\scl});
    \node[font=\scriptsize, text=EHMarginGray, anchor=south west]
      at ({3*\scl*cos(55)}, {3*\scl*sin(55)}) {$r\!=\!3M$};

    % Observer indicator (upper right)
    \draw[->, EHMarginGray, thin]
      (2.6, 1.8) -- ++(0.35, 0.25)
      node[right, font=\scriptsize] {observer};

    % ── n = 0 ray (gold): mild deflection ──────────────────────────────
    \draw[EHAccretionGold, thick, ->]
      plot[smooth, tension=0.55] coordinates {
        (-2.8, -1.0)
        (-1.2, -0.2)
        (0.2, 0.6)
        (1.5, 1.2)
        (2.6, 1.8)
      };
    \node[font=\scriptsize, text=EHAccretionGold]
      at (-2.8, -0.65) {$n = 0$};

    % ── n = 1 ray (blue): half-orbit winding ──────────────────────────
    \draw[EHJetBlue, thick, ->]
      plot[smooth, tension=0.5] coordinates {
        (-2.8, -0.15)
        (-1.6, 0.2)
        (-0.6, 0.9)
        (-0.2, 1.6)
        (0.3, 1.8)
        (1.1, 1.6)
        (1.8, 1.3)
        (2.6, 1.8)
      };
    \node[font=\scriptsize, text=EHJetBlue]
      at (-2.8, 0.2) {$n = 1$};

    % ── n = 2 ray (crimson): full orbit winding ───────────────────────
    % Trajectory hugs the photon sphere (r=3M, radius 1.65cm) more tightly
    \draw[EHHorizonCrimson, thick, ->]
      plot[smooth, tension=0.45] coordinates {
        (-2.8, 0.5)
        (-1.8, 0.8)
        (-1.1, 1.35)
        (-0.2, 1.7)
        (0.6, 1.65)
        (1.3, 1.15)
        (1.55, 0.4)
        (1.3, -0.35)
        (0.6, -0.95)
        (-0.2, -1.15)
        (-0.9, -0.85)
        (-1.2, -0.2)
        (-0.9, 0.5)
        (-0.2, 1.0)
        (0.7, 1.2)
        (1.6, 1.5)
        (2.6, 1.8)
      };
    \node[font=\scriptsize, text=EHHorizonCrimson]
      at (-2.8, 0.85) {$n = 2$};

    % Panel label
    \node[font=\scriptsize, anchor=north] at (0, -2.4) {(a) Ray trajectories};

  \end{scope}

  % ── Panel (b): Observer's image plane ────────────────────────────────

  \begin{scope}[shift={(5.8, 0)}]

    % Shadow: grey filled circle (radius = b_c on image plane)
    % Use 1.2 cm as the shadow angular radius on the image plane
    \def\bshadow{1.2}

    % Shadow fill
    \fill[EHMarginGray, opacity=0.12] (0,0) circle (\bshadow);
    \draw[EHMarginGray, thin] (0,0) circle (\bshadow);

    % n = 0 ring: wide annulus just outside shadow
    % Width ~ 0.55 cm (visually dominant direct image)
    \fill[EHAccretionGold, opacity=0.12]
      (0,0) circle ({\bshadow + 0.55}) -- (0,0) circle (\bshadow);
    \draw[EHAccretionGold, thick] (0,0) circle ({\bshadow + 0.55});
    % Re-draw inner boundary in gold
    \draw[EHAccretionGold, thick] (0,0) circle ({\bshadow + 0.01});

    % n = 1 ring: narrow annulus
    % Width ~ 0.55/23 ≈ 0.024 → exaggerate to 0.10 for visibility
    \def\rone{{\bshadow + 0.01}}
    \draw[EHJetBlue, thick] (0,0) circle ({\bshadow - 0.005});
    \fill[EHJetBlue, opacity=0.20]
      (0,0) circle (\bshadow) -- (0,0) circle ({\bshadow - 0.10});
    \draw[EHJetBlue, thick] (0,0) circle ({\bshadow - 0.10});

    % n = 2 ring: very narrow annulus (barely visible)
    \fill[EHHorizonCrimson, opacity=0.25]
      (0,0) circle ({\bshadow - 0.10}) -- (0,0) circle ({\bshadow - 0.14});
    \draw[EHHorizonCrimson, thick] (0,0) circle ({\bshadow - 0.14});

    % Labels for rings
    \draw[EHAccretionGold, thin, ->]
      ({(\bshadow + 0.28)*cos(35)}, {(\bshadow + 0.28)*sin(35)})
      -- ++(0.7, 0.45)
      node[right, font=\scriptsize, text=EHAccretionGold] {$n = 0$};

    \draw[EHJetBlue, thin, ->]
      ({(\bshadow - 0.05)*cos(145)}, {(\bshadow - 0.05)*sin(145)})
      -- ++(-0.6, 0.45)
      node[left, font=\scriptsize, text=EHJetBlue] {$n = 1$};

    \draw[EHHorizonCrimson, thin, ->]
      ({(\bshadow - 0.12)*cos(215)}, {(\bshadow - 0.12)*sin(215)})
      -- ++(-0.5, -0.5)
      node[left, font=\scriptsize, text=EHHorizonCrimson] {$n = 2$};

    % Shadow label
    \node[font=\scriptsize, text=EHMarginGray] at (0, -0.15)
      {shadow};
    \node[font=\scriptsize, text=EHMarginGray] at (0, -0.45)
      {$b = b_c$};

    % Demagnification annotation
    \draw[EHMarginGray, thin, <->]
      ({\bshadow + 0.55 + 0.08}, -0.9) -- ({\bshadow + 0.08}, -0.9);
    \node[font=\scriptsize, text=EHMarginGray, anchor=west]
      at ({\bshadow + 0.60}, -0.9) {$\delta b_0$};

    \draw[EHMarginGray, thin, <->]
      ({\bshadow + 0.08}, -1.25) -- ({\bshadow - 0.10 - 0.08}, -1.25);
    \node[font=\scriptsize, text=EHMarginGray, anchor=west]
      at ({\bshadow + 0.12}, -1.25) {$\delta b_1 \approx \delta b_0/23$};

    % Panel label
    \node[font=\scriptsize, anchor=north] at (0, -2.4) {(b) Image plane};

  \end{scope}

\end{tikzpicture}
